% !PRINTED 146
\weberpage{146}
%<<CONT
and applies the same permutations to $\mathrm{sn}(u \pm v)$, $\mathrm{cn}(u \pm v)$, $\mathrm{dn}(u \pm v)$ as well. In this way, e.\,g., the twelve formulae (1) to (12) follow from the three formulae (16).

If these permutations are applied to the relations (18), (19), (20) themselves, there results:
\weq{22}{\begin{aligned}
\mathrm{sn}(u+2K) &= -\mathrm{sn}\,u, & \mathrm{sn}(u+2iK') &= \mathrm{sn}\,u,\\
\mathrm{cn}(u+2K) &= -\mathrm{cn}\,u, & \mathrm{cn}(u+2iK') &= -\mathrm{cn}\,u,\\
\mathrm{dn}(u+2K) &= \mathrm{dn}\,u, & \mathrm{dn}(u+2iK') &= -\mathrm{dn}\,u.
\end{aligned}}
The common periods of these three functions are therefore $4K$, $4iK'$; besides these, however, $\mathrm{sn}\,u$ has the period $2iK'$, $\mathrm{dn}\,u$ the period $2K$, $\mathrm{cn}\,u$ the period $2K + 2iK'$.

We give the theorem on $\Theta$-functions (\S\,21) only another expression, by stating the following theorem:

\gesperrt{Every doubly periodic function of $v$ with the periods $2K$, $2iK'$ (which everywhere has the character of a rational function) is, if it is an even function, a rational function of $\mathrm{sn}^2v$, and if it is an odd function, the product of $\mathrm{sn}\,v\,\mathrm{cn}\,v\,\mathrm{dn}\,v$ with a rational function of $\mathrm{sn}^2v$.}

In this way every theorem that refers to homogeneous functions of order zero of $\vartheta$-functions can be converted into a theorem on elliptic functions, and from this one obtains in turn the corresponding theorems on elliptic integrals. One sees at once that the addition formulae (16) are no others than the formulae (17) of \S\,13 in the first section. In the same way the transformation theory of the $\vartheta$-functions furnishes us a solution of the \gesperrt{Jacobi} transformation problem of the elliptic integrals, as we have set it out in \S\,8.

We will verify this for the two principal transformations of the second order, which in \S\,32, as also in \S\,9, we have designated as the \gesperrt{Gauss} and the \gesperrt{Landen} transformation. According to \S\,32, for the \gesperrt{Gauss} transformation
\weq{23}{\frac{\vartheta_{10}\left(0,\dfrac{\omega}{2}\right)\vartheta_{11}\left(u,\dfrac{\omega}{2}\right)}{\vartheta_{00}\left(0,\dfrac{\omega}{2}\right)\vartheta_{01}\left(u,\dfrac{\omega}{2}\right)} = \frac{2\,\vartheta_{01}(u)\,\vartheta_{11}(u)}{\vartheta_{01}^2(u) + \vartheta_{11}^2(u)},}
\weq{24}{\frac{\vartheta_{10}\left(0,\dfrac{\omega}{2}\right)^2}{\vartheta_{00}\left(0,\dfrac{\omega}{2}\right)^2} = \frac{2\,\vartheta_{10}\,\vartheta_{00}}{\vartheta_{00}^2 + \vartheta_{10}^2} = \frac{2\sqrt{\varkappa}}{1+\varkappa},}
%CONT>>
